\section{Completion of the construction}\label{sec:assembly}

We now choose one finite pair of graphs and apply the deterministic results.
\begin{proof}[Proof of Theorem~\ref{thm:main}]
Use the prescribed types of Section~\ref{sec:random} and let
$m\equiv1\pmod4$ tend to infinity. The probability space conditioned on girth
at least $L=\lfloor c_0\log(vm)\rfloor$ is nonempty for sufficiently large
$m$. By Theorem~\ref{thm:no-arrangements}, with probability tending to one
it contains no reduced spherical arrangement. Choose such an outcome with
$L\ge3$, so that both labeled graphs have no loops or parallel edges.

Attach the cones to the rose as in Section~\ref{sec:algebra}, obtaining the
finite two-dimensional complex $X$, and put $G=\pi_1(X)$. By
Proposition~\ref{prop:cone-criterion}, the universal cover of $X$ is
contractible, $G$ is torsion-free, and the labels of paths from either root to
a different vertex of its component are nontrivial in $G$. Finiteness of $X$
gives a finite presentation of $G$.

The intersection parities proved in Proposition~\ref{prop:types} now satisfy
the hypotheses of Proposition~\ref{prop:parity}. That proposition supplies
$a,b,c\in\Ftwo[G]$ with $ab=1$, $ac=0$, and $c\ne0$. Associativity excludes
$ba=1$, as asserted.
\end{proof}

The argument selects a finite matching outcome rather than listing its edges.
Once selected, its cones give a finite presentation, and the sums in
Proposition~\ref{prop:parity} give finite scalar witnesses.
